\cotefile{10}
\scanpage{17}
\margin{\unsure{\ensuremath{\Vert}}}telle que les formes $\varphi$ sur $M$ qui sont des
$(zP)$ polarisations, sont celles pour lesquelles $\varphi^z$
(défini par $\varphi^z(x,y)=\varphi(x,zy)$) sont \deleted{définies} des
$P$-polarisations.

19) De façon générale, \deleted{\unsure{sont}} pour une $\otimes$-catégorie
$\mathcal M$ sur $K\subset\mathbb R$, si $\varphi$ est une forme
$M\otimes M\unsure{\mathord{-}|\longrightarrow}L_{\unsure{\mathrm{Ligne.}|\mathrm{ligne.}}}$, \deleted{\unsure{alors}} et $z\in {}_2\mathrm{Aut}_{\otimes}(\mathrm{id}_{\mathcal M})={}_2Z(K)$
($Z=$ Centre du groupe de \unsure{Galois|choix} \deleted{\unsure{quelconque}} de $\mathcal M$),
alors on définit $\varphi^z$ sur $M\times M$ par
$\varphi^z(x,y)=\varphi(x,zy)$. On trouve alors, si $\varphi^z$ et $\psi$ (sur $N$)
sont non dég. et $\epsilon$-sym ($\epsilon=\pm1$), et si
% slanted
\margin{La condition $z^2=1$ sert \unsure{à|:} assurer que $\varphi$ $\epsilon$-sym implique $\varphi^z$ $\epsilon$-sym.}
\deleted{\unsure{g}} $u:M\longrightarrow N$, $u^{*(\varphi,\psi)}:N\longrightarrow M$, alors
$u^{*(\varphi^z,\psi^z)}=z_M^{-1}u^{*(\varphi,\psi)}z_N=u^{*\varphi,\psi}$. Donc : \deleted{$\varphi$ et $\psi$ sont}

\[
\begin{aligned}
\varphi\ \text{de Weil}&\Longleftrightarrow\varphi^z\ \text{de Weil}\\
\varphi\ \text{et}\ \psi\ \text{de Weil comp.}&\Longleftrightarrow\text{idem pour }\varphi^{\unsure{\epsilon|z}}\text{ et }\psi^\epsilon.
\end{aligned}
\]

D'ailleurs on a $\varphi^\epsilon\otimes\psi^\epsilon=(\varphi\otimes\psi)^\epsilon$. \deleted{\unsure{Donc on}}
\deleted{On en conclut}
\deleted{bien que} Si $z$ opère par id sur $K(1)$
(i.e. $z\in Z'(K)$) alors on voit que si $P$ est
une polarisation, alors les $z(P(M))=\{z\varphi\mid\varphi\in P(M)\}$
($M\in\mathrm{Ob}\,\mathcal M^i$; $i\in\mathbb Z$) forment une polarisation $\deleted{\unsure{z}}\unsure{\epsilon|z}P$ sur $\mathcal M$.
On voit tout de suite que \deleted{cette} dernière
n'est égale à $P$ que si $z=1$. Ainsi,
\margin{? \unsure{!!|\ensuremath{\Vert}}}
$ {}_2Z(\mathbb R)$ \uline{opère librement sur} $\mathrm{Pol}(\mathcal M)$. \uuline{Question}
$\mathrm{Pol}(\mathcal M)$ \unsure{est-il|est il} un pseudotorseur sous $ {}_2Z(\mathbb R)$ ?
% archive sheet number: 16

\scanpage{19}
\margin{\unsure{\ensuremath{\Vert\;\Vert}}}\margin{Notes anciennes}
$\varphi$ forme symétrique resp. alternée sur un motif $M$ de
poids $p$, $\varphi$ à valeurs dans $Q(p)$, $\varphi$ non
degénérée

$\varphi$ définit un isomorphisme
\[
\mathrm{End}(M)\mathrel{\mathop{\longleftarrow}\limits^{\sim}}\deleted{\unsure{\mathrm{Hom}(M,M^{\vee})}}\Phi(M)
\]
en associant à toute forme $\psi$ l'endomorphisme
$\widetilde\psi^{-1}\widetilde\varphi$ de $M$, et à tout
endomorphisme $u$ de $M$ la
forme $\varphi\circ(u\otimes\mathrm{id}$, \deleted{donc} \unsure{correspondant|correspondant à}

\diagram{\begin{tikzcd}[ampersand replacement=\&]
M \arrow[r,shift left=1.4ex,"\widetilde\varphi"] \arrow[r,shift right=1.4ex,"\widetilde\psi"'] \& M^*=M^{\vee}(p)
\end{tikzcd}}

\margin{$\varphi(x,y)=\deleted{\unsure{\varphi}}(y,\widetilde\varphi(x))$}

$\widetilde\varphi\circ u$. Alors la \uline{symétrie} de $\Phi(M)$ devient
une application de $\mathrm{End}(M)$ en lui-$=$, notée $u\leadsto$
$u^\varphi$, qui n'est autre que l'\uline{adjoint}
relativement à $\varphi$,
% archive sheet number: 18

\scanpage{20}
% The two lines run upside down with respect to the archive number.
\deleted{\unsure{hier}. Il \unsure{paraît|faut} \unsure{sans} rapport \unsure{avec|par} \unsure{tes|ta} \unsure{rêveries|relatives|lettres},
\unsure{qui|qu'} \unsure{sont} \unsure{après}}
% archive sheet number: 19

\scanpage{21}
\[
\Phi(M)-\mathcal B(M,M;Q(p))\xrightarrow{\sim}\mathrm{Hom}(M,M^*)\qquad\text{où }M^*=M^\vee(p)
\]
par
\[
\left\langle\underset{\substack{\rotatebox{90}{$\in$}\\ M^*}}{\widetilde\varphi(x)},\underset{\substack{\rotatebox{90}{$\in$}\\ M}}{y}\right\rangle\simeq\varphi(x,y)
\]

Si $\widetilde\varphi$ est un isom, on trouve
\[
\deleted{\unsure{\mathrm{Hom}}}\mathrm{End}(M)\xrightarrow{\overset{\Theta_\varphi}{\sim}}\deleted{\unsure{\mathcal B(M,}}\Phi(M)\simeq\mathrm{Hom}(M,M^*)
\]
en associant à $u$
\[
\Theta_\varphi(u)=\varphi\circ(u\otimes1)\qquad\text{donc }\widetilde{\Theta_\varphi(u)}=\widetilde\varphi\circ u
\]
et en sens inverse
\[
\Theta_\varphi^{-1}(\psi)=\widetilde\varphi^{-1}\widetilde\psi
\]
\diagram{\begin{tikzcd}[ampersand replacement=\&]
 M\arrow[r,shift left=1.4ex,"\widetilde\varphi"]\arrow[r,shift right=1.4ex,"\widetilde\psi"'] \& M^*
\end{tikzcd}}

La symétrie \inserted{$\psi\mapsto\psi'=(-1)^p\,{}^s\psi$} dans $\Phi(M)$ devient un
endomorphisme de $\mathrm{End}(M)$, \deleted{\unsure{par}} savoir $u\mapsto u^\varphi$
\[
u^\varphi=\deleted{\unsure{{}^s\varphi\circ(u\otimes1)}}(-1)^p\widetilde\varphi^{-1}\widetilde{{}^s[\varphi\circ(u\otimes1)]}=\widetilde\varphi^{-1}\bigl(\varphi^s\circ(1\otimes u)\bigr)
\]
de sorte que
\[
\begin{aligned}
&\deleted{\varphi(ux,y)=\varphi(x,u^\varphi y)}\\
&\varphi\unsure{(}u^\varphi x,y)=(-1)^p\varphi(uy;x)
\end{aligned}
\]
\[
(-1)^p\,{}^s\Theta_\varphi(u)^{\prime\prime}(x,y)=\Theta_\varphi(u)^{\prime}(y,x)
\]
donc si $\varphi$ est symétrique \inserted{et $p$ pair}, on trouve
\[
\varphi(ux,y)=\varphi(x,u^\varphi y)\tag*{\deleted{\unsure{($\alpha$)|(a)}}}
\]
si $\varphi$ est \deleted{\unsure{alternée}} \deleted{\unsure{en antisym|ou antisym}}étrique \deleted{\unsure{si $p$ est|et $p$ est}} \inserted{et $p$ impair}, \unsure{églt}
\[
\varphi(ux,y)=\deleted{\unsure{(-1)}}\varphi(x,u^\varphi y)\tag*{\unsure{($\alpha$)|(a)}}
\]
% archive sheet number: 20

\scanpage{22}
\deleted{\unsure{bien|rien|tion}, c'est pourquoi il faut surtout maintenant}
% archive sheet number: 21

\scanpage{23}
Les éléments fixes sous l'involution $u\mapsto u^\varphi$ correspondent
alors aux formes \unsure{symm.|sym.} resp. alternées
(suivant $p$ pair ou impair). De plus, on a
\[
\left\{\begin{aligned}
(u+v)^\varphi&=u^\varphi+v^\varphi\\
(\lambda u)^\varphi&=\lambda u^\varphi && \lambda\in Q\\
(u^\varphi)^\varphi&=u\\
(uv)^\varphi&=\deleted{\unsure{\varphi}}v^\varphi u^\varphi
\end{aligned}\right.
\]

Si on remplace $\varphi$ par une forme $\psi$ de $=$
type (pair resp. impair) on aura $\psi\longleftrightarrow a$
\[
\psi(x,y)\deleted{\unsure{(x,\psi)}}=\varphi(ax,y)
\]
et l'hyp. sur $\psi$ s'exprime, comme on a dit
par $a^\varphi=a$. On aura de plus
\[
u^\psi=\deleted{\unsure{a^\varphi}}a^{-1}u^\varphi a\qquad\text{($a$ étant inv.).}
\]
donc l'\uline{involution de} $A=\mathrm{End}(M)$ \uline{est
bien déterminée mod autom intérieur}
(indépendamt du choix de $\varphi$, ou $\psi$), une
fois si $\varphi$ est \unsure{fixée|fixé}, mod aut. int.
par un élément \deleted{de $A$ symétrique} \inserted{$\varphi$-hermitien} de $A$.
% archive sheet number: 22

\scanpage{24}
On notera cependant que la donnée d'une
forme $\varphi$ \inserted{sur $M$} est \uline{plus} \uline{précis} que la donnée de
l'involution sur $A=\mathrm{End}(M)$ ; ainsi si $\mathrm{End}(M)$
est commutatif, ttes les $\varphi$ donnent la
$\deleted{\unsure{m}}=$ involution ! De $=$, si $\lambda\in Q^*$, $\lambda\varphi$ donne
la $=$ involution que $\varphi$

Considérons la forme \uline{symétrique}
\[
\mathrm{Tr}\,v^\varphi u=\mathrm{Tr}\,u^\varphi v
\]
sur \inserted{le motif} \deleted{$A=$} $\underline{\mathrm{End}}(M)\simeq$ \deleted{\unsure{induite}, par la forme} \inserted{$M^\vee\otimes M$}
cette forme n'est autre que $\overset{\vee}{\varphi}\otimes\varphi$ où
$\overset{\vee}{\varphi}$ est la forme \uline{duale} de $\varphi$ sur $M^\vee$, \inserted{déduite de $\varphi$ par transport de structure par $\widetilde\varphi$} \deleted{\unsure{ou}}
\deleted{\unsure{auquel} on \unsure{a} identifié $\underline{\mathrm{End}}(M)$ à
$\unsure{\Phi}(M)\simeq M^\vee\otimes M^\vee$ et \uline{supposons} qu'elle soit
dé}
Elle induit une forme sur
$\mathrm{End}(M)$ à valeurs ds $Q$, \unsure{notre point|notre forme|motif point|notée point}.

\uline{Supposons} que cette forme soit \deleted{\uline{définie}}
\uline{positive} (elle est néc. non dég, donc \uline{définie}
positive). \deleted{Soit $P$ le cône des} Alors
$A=\mathrm{End}(M)$ devient une ``algèbre de Riemann''.
% archive sheet number: 23

\scanpage{25}
Soit $A^+$ son cône positif, \deleted{engendré} adhérence
de l'ens des combinaisons linéaires à
coefficients $\geq0$ des $u^\varphi u$ ($u\in A$ ; on
peut $=$ supposer $u\in A^*$, groupe des \uline{unités} de $A$)

Il correspond dans $\Phi^s$ au\deleted{\unsure{x}} cône fermé $\Phi^+$
engendré par les $\varphi^u$ $[\varphi^u(x,y)=\varphi(ux,uy)]$,
i.e. le plus petit cône fermé dans $\Phi^s$
contenant $\varphi$ et stable par les automorphismes
de $M$. \deleted{Tout élément de} De la théorie des
Algèbres de Riemann résulte que tt
élément $a$ \uline{intérieur} à $A^+$ est une
unité, donc tt élément intérieur à $\Phi^+$
est une forme non dégénérée ; \deleted{\unsure{Donc on}} \inserted{de plus,}
pour un tel $a$, \deleted{\unsure{l'involut|l'involution}} (correspondant à
\deleted{telle que \unsure{engendre} \unsure{le|une des} $=$ cône}
$\varphi\in\Phi^+$ intérieur à $\Phi^+$) l'involution
$u\mapsto u^\psi=a^{-1}u^\varphi a$ \deleted{satisfait à} étant déduite
\[
\deleted{\unsure{\mathrm{Tr}}\,u^\psi u=\unsure{\mathrm{Tr}}\,a^{-1}u^\varphi au=\sum_i b_i^{\deleted{\unsure{*}}-1}(b_i^\varphi)^{-1}u^\varphi b_i^\varphi b_i u}
\]
de la précédente par un autom (intérieur)
qui ne bouge pas la trace, elle \deleted{\unsure{peut}} possède
% archive sheet number: 24

\scanpage{26}
la $=$ propriété de positivité, savoir
\[
\mathrm{Tr}\,u^\psi u\deleted{\unsure{=}}\quad(=\mathrm{Tr}\,(a^{-1}ua)^\varphi(a^{-1}ua))\geq0
\]
égalité seulement si $u=0$. Donc tt pt de
l'intérieur du cône \deleted{de} $\Phi^+$ engendre
(via autom de $M$) le $=$ cône $\Phi^+$.

C'est la donnée du cône $\Phi^+$ qui,
intrinsèquement, constitue une structure
supplémentaire importante.
% slanted
\margin{dégager la structure abstraite : anneau $A$ \inserted{de type fini sur $Q$}, bimodule $\Phi$ sur $A$ libre à gauche et à droite, Trace sur $A$ ; notion de \uline{bon} cône ds $\Phi$}

On notera que, sous cette forme abstraite,
on ne peut (parmi les cônes possibles,
correspondant aux divers choix de $\varphi$)
en distinguer de meilleurs que d'autres.
Ainsi, si $\varphi$ satisfait la condition de
positivité de départ, de $=$ $-\varphi$, qui
définira le cône $-\Phi^+$

Noter aussi : \deleted{Tout cône convexe fermé
qui contient} si $v,w\in A^+$,
alors $\mathrm{Tr}\,vw\geq0$ [et
en fait, $A^+$ est son \uline{propre} \uline{polaire} dans
$A^s$ \inserted{hermitiens} pour la forme $\mathrm{Tr}\,\unsure{uv|vw}$]. Or pour
% archive sheet number: 25

\scanpage{27}
\deleted{$\Theta_\varphi$, la forme $\mathrm{Tr}\,uv$ devient
la forme sur $\Phi=\unsure{\Gamma|\mathscr T}(M^\vee\otimes M^\vee(p))$ \unsure{de} induite
par $\overset{\vee}{\varphi}\otimes\overset{\vee}{\varphi}(p)$, ou encore,}
identifiant $A$ d'une part à
$\Phi(M)=M^\vee\otimes M^\vee(p)$ \inserted{par $\Theta_\varphi$}, d'autre part via
\deleted{\unsure{$\varphi$}} $u\mapsto{}^tu$ et $\overset{\vee}{\varphi}$ à [\unsure{ou|où} via $\Phi(M)\simeq\Phi(M^{*})$
à l'aide de l'isom $\widetilde\varphi:M\longrightarrow M^*$] à
$\Phi(M^*)=M\otimes M(-p)$, la forme $\mathrm{Tr}\,uv$
sur $A\times A$ devient la forme évidente
de contraction sur $\Phi(M)\times\Phi(M^*)$.

Noter d'autre part que \deleted{pour} \inserted{\unsure{à}} tt ``cône
\deleted{positif sur} admissible'' $\Phi^+(M)$ de $\Phi(M)$,
est canoniquement associé un ``cône admissible''
$\Phi^+(M^*)$ de $\Phi(M^*)$, en prenant les
$\overset{\vee}{\varphi}^*$ ($\varphi\in\Phi^{\unsure{+|*}}(M^*)^\circ$). Ceci posé, on voit
que \uline{$\Phi^+(M)$ et $\Phi^+(M^*)$} sont
polaires l'un de l'autre. Conséquence :
% archive sheet number: 26

\scanpage{28}
Si on a deux cônes \inserted{convexes} \unsure{fermé|fermés} $P,P^*$ dans
$\Phi(M)$, $\Phi(M^*)$, \deleted{qui sont} tels que

a) $P,P^*$ stables par automorphismes

b) $\langle P,P^*\rangle\geq0$

c) Il existe $\varphi,\varphi^*$ polaires l'un de l'autre,
``admissibles'' [i.e. \deleted{\unsure{$\mathrm{Tr}\,u$|Pour}} $\varphi$ non dég. et
$\mathrm{Tr}\,u^\varphi u\geq0$ pour tt $u\in A$], avec $\varphi\in P$, $\varphi^*\in P^*$

alors $P,P^*$ sont néc. deux cônes admissibles
associés, (et en particulier chacun est
exactement le polaire de l'autre, et $\overset{\circ}{P}^{\deleted{\unsure{*}}+}$
consiste exactement en les $\overset{\vee}{\varphi}$ pr $\varphi\in\overset{\circ}{P}^+$).

\uline{Abstraitement}

Nous avons une Algèbre $A$ $\frac12$ simple de type fini sur $Q$,
\inserted{munie d'une forme $\mathrm{tr}$ telle que $\mathrm{tr}(uv)$ \uline{\unsure{non dégénérée|non définie}}}
un \deleted{\unsure{b|bi}}module \inserted{à droite} $\Phi$ sous $A\otimes A$ (élément $\varphi$,
opération de $A$ écrite $u\varphi v$) \deleted{\unsure{biunivoque|inversible|transportés}} libre
de rang $1$ sur $A\otimes1$ et sur $1\otimes A$
\deleted{(\unsure{donc canoniqt|isom. correspond.} $\unsure{\simeq}$ à $A$)}

le \deleted{bimodule} $Q$-dual $\Phi^*$, \deleted{\unsure{églt et isomorphe}} muni de sa
structure de module sur $A^\circ\otimes A^\circ$
\deleted{à $A$ par conséquent}

Pour un élément $\varphi$ de $\Phi$, il revient au $=$
de dire qu'il est une base \deleted{\unsure{à gauche}
à droite, \unsure{Je lui} \unsure{fais} \unsure{on|se} dit \unsure{qu'il est} non dég.}
% archive sheet number: 27

\scanpage{29}
\deleted{Il lui correspond alors $\varphi^{-1}$ un élément
\deleted{unique}\inserted{canonique}
de $\Phi$ \unsure{il forme donc}}

sur $1\otimes A$ ou sur $A\otimes1$. L'utilisant
\deleted{pour} écrire ces \unsure{deux|deux iso|isom.} de modules sur
$A\otimes1$ \unsure{$\Phi\rightleftarrows A$|$\Phi\simeq A$}, et \deleted{utilisant}\inserted{composant avec} $\mathrm{tr}$ sur $A$
\deleted{pour écrire $A^*\simeq\mathrm{dual}_Q$ de $A\simeq A$}, on
trouve une forme sur $\Phi$, i.e. un élément noté
$\check\varphi$. Procédant à droite au lieu de prendre
à gauche on trouve pareil.

On \deleted{\unsure{veut|peut}} utilise une \uline{symétrie} sur $\Phi$ qui
échange les deux structures de $A$-modules.
L'application $\overset{\circ}{\Phi}{}^{\deleted{\unsure{*}}}\longrightarrow\overset{\circ}{\Phi}^{*}$
commute néc. à cette symétrie.

Pour tt élément $\varphi$ de $\overset{\circ}{\Phi}^s$ [fixe par
symétrie] on trouve une bijection
$A\longrightarrow\Phi$, d'où sur $A$ une \deleted{involution}
application $u\leadsto u^\varphi$ linéaire, telle que
$(u^\varphi)^\varphi=u$, [déterminée mod \deleted{conjugaison}
\inserted{\uuline{composition} \deleted{\unsure{avec}} aut. int.}
par un élément $a$ de $A$ tel que $a^\varphi=a$, si
on change $\varphi$ en $\psi$].
\margin{\tikz[baseline=(a.base)]{\node[inner sep=2pt] (a) {\shortstack[l]{$\unsure{Z|2}$\\\unsure{$\mathord{/}_{\circ}$|6|$\delta$} (mais non\,: conjugaison\\ par !)}};\draw (a.center) ellipse [x radius=1.6cm,y radius=.9cm];}}
Si on fait la $=$ construction pour $\check\varphi$, on trouve une
involution sur $A^\circ=A$ (ensemblistement)
qui est la \deleted{\unsure{sa}}$\mathord{/}^{\circ}$.
% slanted
\margin{$\mathord{/}^{\circ}$ On a ici $(uv)^\varphi=v^\varphi u^\varphi$ (indépendant
des deux particulier de $\psi$)}
La structure de \inserted{\unsure{départ}} bimodule de
$\Phi$ devient la structure $aub^\varphi$, et $\varphi\leadsto\varphi^{-1}$ devient
\inserted{$u\leadsto u^{-1}$} \ldots
% archive sheet number: 28

\scanpage{30}
On dit que \deleted{$(A,\Phi)$} \deleted{\unsure{$\varphi$|$\Phi$}} $\varphi$ est \deleted{\unsure{$\varphi$}} \uline{\unsure{riemannien|riemannienne}} si $A$,
muni de sa trace et de l'involution $u\leadsto u^\varphi$, est
une Algèbre riemannienne, i.e. satisfait
\[
\mathrm{Tr}\,u^\varphi u\geq0\qquad \text{tt }u
\]
\deleted{ce qui ne dépend pas du choix de $\varphi$}

[\uline{NB} utilisant les identifications $A\simeq\Phi$, $A^\circ\simeq\Phi^*$
définies par $\varphi$ et $\varphi^{-1}$, la forme $\mathrm{Tr}\,v^\varphi u$
sur $A^\circ\times A$ s'identifie à la forme canonique
sur $\Phi\times\Phi^*$, elle est donc \uline{non dégénérée}]

Tout $\varphi$ riemannienne définit comme plus
haut un cône convexe fermé $\Phi^+(\varphi)=\Phi^+$,
savoir le plus petit cône convexe fermé
stable par $u\otimes u$ ($u\in A$) contenant $\varphi$, et qui
correspond par $\Theta_\varphi:A\simeq\unsure{\Phi|P}$ au cône des
``éléments positifs'' de l'algèbre \deleted{hilb.} riemannienne
$A$, engendré par les éléments de la forme
$u^\varphi u$. De $=$, $A^\circ\longrightarrow\Phi^*$ transforme ce dernier
cône en le cône $\Phi^{*+}$ défini en termes
de $\varphi^{-1}$. La théorie des Alg. hilbertiennes
nous dit que ces cônes $\Phi^{\deleted{\unsure{*}}+}$ et $\unsure{\Phi|\varphi|P}^{*+}$ sont
polaires l'un de l'autre, \unsure{ts.|tq.} que leurs
% archive sheet number: 29

\scanpage{31}
intérieurs correspondent aux formes dans les
cônes qui sont \uline{non dégénérées}, ou encore
(étant toutes \inserted{pré}riemanniennes \unsure{églt|égalmt}) riemanniennes,
et que pour tt $\psi$ de l'intérieur de
$\Phi^+(\varphi)$ le cône $\Phi^+(\psi)$ est égal à $\Phi^+(\varphi)$.

On trouve ainsi dans $\Phi$ un certain
nombre de cônes convexes \inserted{fermés} $P_\alpha$ avec
intérieur, stables par les $u\otimes u$ ($u\in A$)
et dont les intérieurs sont deux à
deux disjoints, \deleted{\unsure{leur}} \unsure{la|leur} réunion \inserted{des intérieurs} étant l'ens
des formes riemanniennes relativement à $\Phi$.

S'il y a un $\varphi$ riemannienne, alors $-\varphi$ est
riemannienne, mais certainement pas dans le
$=$ cône. (les deux cônes sont opposés).

Quel est la signification de l'invariant :
``nb des cônes riemanniens dans $\Phi$''
(qu'on peut aussi considérer comme
un invariant pour une Alg. riemannienne
sur $Q$, ou seulement sur $\mathbb R$ \unsure{---|-,-}].
% archive sheet number: 30

\scanpage{88}
Filtration \unsure{de|des} foncteurs fibres
% archive sheet number: 87
